\documentclass[10pt,a4paper]{article}

\usepackage[T1]{fontenc}
\usepackage[utf8]{inputenc}
\usepackage{lmodern}
\usepackage{amsmath,amssymb,amsthm}
\usepackage{booktabs,array,tabularx}
\usepackage{geometry}
\usepackage{xcolor}
\usepackage{enumitem}
\usepackage{microtype}
\usepackage{seqsplit}
\usepackage{hyperref}
\usepackage{fancyhdr}

\geometry{margin=1in}
\hypersetup{
  colorlinks=true,
  linkcolor=blue!55!black,
  citecolor=blue!55!black,
  urlcolor=blue!55!black,
  pdftitle={The Jevil Is in the Details: State Recovery and Fresh-Message Forgery Below Jevil's Signature Cliff},
  pdfauthor={Martin Feussner},
  pdfsubject={Cryptanalytic audit report for IACR ePrint 2026/1103},
  pdfkeywords={digital signatures, Jevil, zk-WHIR, Reed--Solomon, state recovery, forgery}
}
\setlist{nosep,leftmargin=*}
\setlength{\parindent}{1em}
\setlength{\parskip}{0.25em}
\setlength{\emergencystretch}{2em}
\pagestyle{fancy}
\fancyhf{}
\fancyhead[L]{The Jevil Is in the Details}
\fancyhead[R]{15 October 2026}
\fancyfoot[C]{\thepage}
\renewcommand{\headrulewidth}{0.3pt}

\newtheorem{theorem}{Theorem}
\newtheorem{lemma}{Lemma}
\newtheorem{corollary}{Corollary}
\newcommand{\F}{\mathbb F}
\newcommand{\RS}{\mathsf{RS}}
\newcommand{\Commit}{\mathsf{Commit}}
\newcommand{\Open}{\mathsf{Open}}
\newcommand{\Verify}{\mathsf{Verify}}
\newcommand{\Sign}{\mathsf{Sign}}
\newcommand{\KeyGen}{\mathsf{KeyGen}}
\newcolumntype{Y}{>{\raggedright\arraybackslash}X}
\newcommand{\hash}[1]{{\ttfamily\small\seqsplit{#1}}}

\title{\bfseries The Jevil Is in the Details:\\
State Recovery and Fresh-Message Forgery Below Jevil's Signature Cliff}
\author{Martin Feussner\\
\small Selmer Center, University of Bergen\\
\small \texttt{martin.feussner@uib.no}}
\date{15 October 2026}

\begin{document}
\maketitle

\noindent\fcolorbox{black!50}{black!4}{\parbox{0.94\textwidth}{
\textbf{AI disclosure and review status.}
This audit used the OpenAI model Daybreak Blue at the ultra reasoning level.
OpenAI Codex performed AI-assisted target discovery, attack-history searching,
scheme and source review, cryptanalysis, implementation, experimentation, and
technical writing. Specialist AI agents worked on bounded components of the
analysis and implementation. A separately tasked hostile AI reviewer replayed
and challenged the primary result, and a fresh-context AI agent independently
implemented and reproduced the attack on a new key. These were AI review and
AI reproduction, not human verification. The primary experiment, hostile
adversarial review, fresh-context independent reproduction, final exact-target
attack-history review, and independent Publication Review passed. No independent
human verification was performed; it is optional for this audit. OpenAI Codex
generated and selected the title and intentionally chose ``The Jevil Is in the
Details'' as scheme-related wordplay on ``the devil is in the details,'' because
the flaw lies in Jevil's width-four opening and encoding details.}}

\begin{abstract}
Jevil is a post-quantum few-time signature scheme designed to remain secure for
$n^\star=1023$ ordinary signatures and to reveal its signing polynomial at the
advertised 1,024th-signature cliff. We give a different, practical attack below
that cliff against the current fourth-revision, source-backed profile. Each
ordinary zk-WHIR proof authenticates 64 positions of a reused initial
commitment, and each position discloses a width-four Reed--Solomon leaf. At this
profile the committed object has four rows of degree below 32,768. Consequently,
32,768 distinct leaf positions determine every row, including the polynomial
coefficients and commitment randomness needed to rebuild an equivalent signing
state.

Our primary public-data process used 641 ordinary signing queries, authenticated
41,024 openings at 35,203 distinct positions, recovered all four rows, rebuilt
the exact public root and evaluation, and produced a signature on an unqueried
message. The unmodified ordinary verifier accepted it. A hostile reviewer
replayed all public inputs and produced a second accepted signature. A
fresh-context implementation independently repeated the full attack on a new
key with 650 queries and 35,603 distinct positions; a second untouched-source
verifier accepted the forgery. Changed-message, changed-root, corrupted-row,
and altered-signature controls rejected.

The executable claim is pinned to the source commit identified in
Section~\ref{sec:target}. The same full-recovery threshold also follows
analytically for the intended four-row encoding of zk-WHIR
Definition~3.22. We keep that analytic scope distinct from the source's
nonconforming contiguous row allocation. We make no claim that Reed--Solomon
interpolation is new, no minimum-query claim, and no broad priority claim about
all Jevil development versions.
\end{abstract}

\section{Introduction}

Jevil is a transparent post-quantum few-time signature whose defining feature
is a sharp, count-triggered failure point~\cite{jevil}. For the largest
source-backed row, a signer may issue at most $n^\star=1023$ signatures. Each
signature publishes $K=16$ message-derived evaluations of a secret polynomial
of dimension $M=(n^\star+1)K=16{,}384$. The paper deliberately chooses this
dimension so that signature 1,024 can provide enough evaluations to recover the
polynomial. Below that boundary, the paper claims more than 124 bits of
classical EUF-CMA-$n^\star$ security.

The deliberate cliff counts the 16 explicit message-derived evaluations. The
zk-WHIR proof attached to every signature exposes a second and larger stream of
reused-commitment data: 64 authenticated initial-codeword leaves. Each leaf
contains four field elements. At the audited endpoint those elements are one
evaluation from each of four degree-below-32,768 Reed--Solomon rows. Accumulating
32,768 distinct outer positions therefore recovers the complete committed
state. The opening interface accepts fresh prover randomness separately from
that state, so the recovered state is an equivalent signing witness.

The attack has five concrete steps:
\begin{enumerate}
  \item request ordinary signatures on distinct messages;
  \item verify them and extract only their authenticated initial width-four
        openings;
  \item deduplicate positions and interpolate the four rows after 32,768
        distinct positions have appeared;
  \item re-encode the recovered rows, require the exact public root and public
        polynomial evaluation, and reconstruct signer state; and
  \item use fresh opening randomness to sign an unqueried message and submit
        the result to the unchanged ordinary verifier.
\end{enumerate}

This is ordinary erasure recovery for Reed--Solomon codes, followed by state
reconstruction. Reed--Solomon interpolation and fast multipoint algorithms are
standard~\cite{reedsolomon,mca}; we make no methodological-novelty claim. The
result is the application of that standard machinery to Jevil's accumulated
width-four openings, with a full-parameter public-data implementation and an
ordinary-verifier-accepted fresh-message forgery.

\paragraph{Contributions and scope.}
We identify the unit mismatch between the claimed hiding budget and the
width-four leaves; prove full recovery under the paper's intended
Definition~3.22 encoding and the pinned source encoding; implement the attack
at $n^\star=1023$ without using signature 1,024; and report a primary run,
hostile replay, and independently written fresh-key reproduction. The
executable result applies to one precisely pinned source-backed profile. The
paper/source correspondence and encoding distinctions in
Sections~\ref{sec:target} and~\ref{sec:encodings} are part of the claim.

\section{Exact target and claim boundary}
\label{sec:target}

\subsection{Paper and source pin}

The paper target is Nadim Kobeissi's \emph{Jevil: A
Catastrophic-Failure-by-Design Signature Scheme}, IACR ePrint 2026/1103,
current revision 4~\cite{jevil}. The pinned 404,750-byte PDF has SHA-256 digest
\begin{center}
\hash{c0da79bdf8b2d2c9667d6167d6ab008b8eeaf838b463a40eb64f2abe123c01c5}.
\end{center}
Only revision 4 of the four reported revisions was available for bytewise
inspection.

The executable target is the author's \texttt{symbolicsoft/jevil} repository at
commit
\begin{center}
\hash{905466fbe3be99d90fd41e461e0c68271d09104b}
\end{center}
and Git tree
\begin{center}
\hash{f4a7366dac1a1f079186d87a2aca79c9918ee79e}.
\end{center}
The source archive has SHA-256
\hash{b8817b85a5f90dd4045c8e635dad806e3973b26dbb36c01d7d2a37fcdff0270e}.
The commit predates the final PDF by about 19 hours 47 minutes, and the final
paper does not name a source commit. We therefore call this the
\emph{source-backed current-revision profile}, rather than asserting
byte-exact source conformance to every sentence of the PDF.

\subsection{Frozen parameter row}

Table~\ref{tab:parameters} records the values material to the attack. The
working field is a degree-four extension of the Goldilocks prime field. The
source-backed Poseidon2 capacity is $c=4$. The paper discusses $c=6$ as a
possible higher quantum-collision-security choice, but that row is not the
executed target.

\begin{table}[ht]
\centering
\small
\begin{tabular}{@{}ll@{}}
\toprule
Quantity & Frozen value \\
\midrule
Ordinary signing budget $n^\star$ & 1,023 \\
Advertised cliff signature $n^\star+1$ & 1,024 \\
Message-derived evaluations per signature $K$ & 16 \\
Secret polynomial dimension $M$ / degree bound & 16,384 / $<16{,}384$ \\
HVZK padding $B$ & 114,688 \\
Initial encoded message length $N=M+B$ & 131,072 \\
Main-code rate / width & $1/4$ / 4 \\
Row count / row dimension & 4 / 32,768 \\
Initial outer domain & 131,072 width-four leaves \\
Initial openings per signature $\theta$ & 64 \\
Poseidon2 capacity & 4 \\
Public key / reported signature size & 68 bytes / about 337 KB \\
\bottomrule
\end{tabular}
\caption{The exact $n^\star=1023$ source-backed profile.}
\label{tab:parameters}
\end{table}

The paper claims EUF-CMA-$n^\star$ security above 124 classical bits for at
most $2^{128}$ random-oracle queries. At capacity four, its own analysis places
the quantum collision floor at about 85 bits; that separate generic floor is
not used in our attack.

\subsection{Dated attack-history statement}

No prior public practical signing attack against the current fourth-revision
Jevil profile at $n^\star=1023$, $c=4$, using at most 1,023 ordinary signing
queries was located in the searches performed through 15 October 2026. This is
a bounded, dated search result. It is not a claim that no Jevil development
version was ever attacked. Section~\ref{sec:history} discloses a
pre-publication root-binding forgery fix and a later Merkle-helper report.

\section{Jevil and its two disclosure streams}
\label{sec:scheme}

\subsection{Signing polynomial and advertised cliff}

Let
\[
  f(X)=\sum_{j=0}^{M-1} c_jX^j\in\F[X],
  \qquad c=(c_0,\ldots,c_{M-1})\in\F^M.
\]
Key generation derives $c$ from a 32-byte seed. It commits once to $c$ using
zk-WHIR, deriving commitment padding $r_{\rm zk}\in\F^B$, and publishes
\[
  \mathsf{pk}=(\mathsf{root},w,n^\star),
  \qquad w=f(z),
\]
where $z$ is derived from the root. The commitment state is cached and reused
for every signature under the key.

For a message $m$, the signer derives $K=16$ positions and corresponding field
points $x_1,\ldots,x_K$, publishes $y_t=f(x_t)$, and batches those evaluations
with $(z,w)$ into one linear-form claim. In simplified notation, Fiat--Shamir
challenges $\beta_1,\ldots,\beta_{K+1}$ define
\[
  \alpha=\sum_{t=1}^{K}\beta_t(1,x_t,\ldots,x_t^{M-1})
     +\beta_{K+1}(1,z,\ldots,z^{M-1}),
\]
\[
  v=\sum_{t=1}^{K}\beta_ty_t+\beta_{K+1}w.
\]
The proof $\pi\leftarrow\Open(\mathsf{st},\alpha,v;\rho)$ uses a fresh
32-byte prover-randomness seed $\rho$. The signature is
$((y_1,\ldots,y_K),\pi)$, and ordinary verification re-derives the positions and
checks the proof against $\mathsf{root}$.

Jevil chooses $M=(n^\star+1)K$. Its advertised cliff is based on accumulating
the explicit $(x_t,y_t)$ pairs: signature 1,024 supplies the next bucket of 16
evaluations from which the degree-below-16,384 polynomial can be interpolated.
That 1,024th signature is outside the permitted budget and outside our attack.

\subsection{The initial-codeword openings}

The proof itself publishes another stream. In its first round, the verifier
derives 64 in-domain positions and authenticates the corresponding initial
commitment leaves by Merkle multiproof. Jevil sets the interleaving width to
four, so each authenticated outer leaf exposes four field elements. The same
initial root is reused across all ordinary signatures.

The paper's multi-opening argument charges $\theta=64$ positions per signature
against $B=N-M=114{,}688$ scalar hiding slots. This mixes units. A width-four
outer position supplies four scalar values; in the cited four-row construction,
one position supplies one value from each row, each of which has only $B/4$
random coefficients. The relevant outer-position hiding allowance is therefore
$B/4=28{,}672$, while the paper permits $n^\star\theta=65{,}472$ outer draws.

This accounting defect has a direct recovery consequence. Each row has
dimension $N/4=32{,}768$. Once the accumulated signatures authenticate that
many distinct outer positions, the entire committed state is public.

\section{Encoding interpretations and full-state recovery}
\label{sec:encodings}

The paper and source do not allocate the encoded state identically. Keeping the
three coherent readings separate prevents a source-layout defect from being
mistaken for the design-level cause of full recovery.

\begin{table}[ht]
\centering
\scriptsize
\begin{tabularx}{\textwidth}{@{}p{0.18\textwidth}YYp{0.19\textwidth}@{}}
\toprule
Reading & Encoded object & One outer leaf & Full-recovery threshold \\
\midrule
Literal paper wording
& One scalar $\RS[4N,N]$ codeword packed into disjoint width-four stripes
& Four scalar evaluations
& 32,768 stripes provide $N$ evaluations \\
Definition~3.22
& Four $\RS[N,N/4]$ rows after separately splitting $c$ and $r_{\rm zk}$
& One evaluation from each row
& 32,768 positions interpolate every row \\
Pinned source
& Concatenate $c\,\|\,r_{\rm zk}$, split it into four contiguous chunks, and
  encode each as an $\RS[N,N/4]$ row
& One evaluation from each row
& 32,768 positions interpolate every row \\
\bottomrule
\end{tabularx}
\caption{Three encoding readings. They differ in coefficient allocation but
share the same full-recovery threshold. The literal paper stripe map is not
fully pinned.}
\label{tab:encodings}
\end{table}

In the intended Definition~3.22 allocation~\cite{zkwhir}, each row contains
4,096 coefficients from $c$ and 28,672 coefficients from $r_{\rm zk}$. In the
pinned source, row zero contains all 16,384 coefficients of $c$ followed by
16,384 padding coefficients, while rows one through three contain only
padding. This source layout changes the earliest partial-leakage threshold. It
does not change the degree bound or the 32,768-position threshold for complete
recovery.

\begin{lemma}[Width-four interpolation]
\label{lem:interpolation}
Let $p_0,\ldots,p_3\in\F[X]$ have degree less than $k=32{,}768$. Suppose an
authenticated outer leaf at domain point $a_i$ reveals
$(p_0(a_i),\ldots,p_3(a_i))$. Any $k$ leaves at distinct domain points uniquely
determine all four polynomials.
\end{lemma}

\begin{proof}
For each fixed row $j$, the leaves provide $k$ evaluations
$(a_i,p_j(a_i))$ at distinct points. A polynomial of degree below $k$ is
uniquely determined by those evaluations. Interpolate independently for
$j=0,1,2,3$.
\end{proof}

For the source layout, concatenating the four recovered coefficient vectors
recovers exactly $c\,\|\,r_{\rm zk}$. For the Definition~3.22 layout, the
specified inverse split recovers the same two logical objects. Under the
literal scalar-code reading, 32,768 disjoint stripes provide $N=131{,}072$
scalar evaluations and determine the single degree-below-$N$ encoding
polynomial, provided the intended stripe-to-domain map is the disjoint packing
described in the paper.

\begin{theorem}[Equivalent signing-state recovery]
\label{thm:state}
For either four-row layout in Table~\ref{tab:encodings}, 32,768 distinct,
authenticated initial leaves from one reused Jevil commitment determine an
equivalent signing state. The state can produce fresh ordinary signatures
without recovering the original 32-byte signing seed.
\end{theorem}

\begin{proof}
Lemma~\ref{lem:interpolation} recovers every row. Re-encoding those rows
reproduces the initial codeword and Merkle tree, hence the public root. Their
coefficient allocation reconstructs $c$ and $r_{\rm zk}$, while the public
relation $w=f(z)$ provides an independent binding check on $c$. These values
determine the commitment state used by $\Open$. The opening API takes fresh
randomness $\rho$ separately, so an attacker may sample unrelated new proof
coins and run the ordinary signing computation. Verification depends on the
public root, message, claimed evaluations, and proof; it does not require the
original seed.
\end{proof}

Theorem~\ref{thm:state} concerns complete state recovery, not merely leakage.
It survives the source's row-allocation error. Our executable demonstration
uses the pinned source layout. The Definition~3.22 statement is an analytic
design-level argument; we did not silently present the source execution as a
separate paper-faithful implementation.

\section{Attack algorithm}
\label{sec:algorithm}

The adversary receives a public key $\mathsf{pk}$ and ordinary signing-oracle
access under the advertised $n^\star=1023$ budget.

\begin{enumerate}
  \item Choose distinct messages $m_0,m_1,\ldots$ and request ordinary
        signatures. Never request query index 1,023, which would be signature
        number 1,024.
  \item Run the unchanged ordinary verifier on every response. Replay its
        Fiat--Shamir transcript to derive the 64 first-round positions, verify
        each initial Merkle opening against the public root, and retain the
        four field elements at each position only after the complete proof has
        verified and the transcript has ended.
  \item Deduplicate positions. Abort on conflicting values. Stop collection
        as soon as at least $k=32{,}768$ distinct positions are available.
  \item Select $k$ positions. Fast-interpolate each of the four
        degree-below-$k$ rows, then re-encode every row over the full
        131,072-point domain. Require every retained opening, including those
        not selected for interpolation, to match the re-encoding.
  \item Reconstruct the source's $c\,\|\,r_{\rm zk}$ layout. Require
        $w=f(z)$ and rebuild the commitment, requiring its Merkle root to equal
        the target public root byte for byte.
  \item Build a signer cache from the recovered state, choose a message
        $m^\star$ absent from the oracle query set, and draw fresh proof
        randomness. Run the ordinary signing path on $m^\star$ and submit the
        result to the unmodified ordinary verifier.
\end{enumerate}

The extractor in step 2 observes only data already sent to an ordinary
verifier. The primary and independent recovery processes had no secret-key or
original-cache input. Root reconstruction and $w=f(z)$ make coefficient order,
field basis, domain indexing, and row assembly fail closed except under the
underlying binding/collision assumptions.

\begin{corollary}[Fresh-message forgery]
If the algorithm reaches step 6, it produces an efficient fresh-message
forgery and a reusable equivalent signing capability for the pinned profile.
\end{corollary}

\begin{proof}
By Theorem~\ref{thm:state}, the reconstructed state opens the exact public
commitment. The message is outside the signing-query set and the prover coins
are freshly chosen. Completeness of the ordinary opening and verification paths
therefore gives acceptance. The same recovered state can be reused with new
messages and fresh coins.
\end{proof}

\section{Sample complexity and success probability}
\label{sec:probability}

In the Fiat--Shamir random-oracle model, distinct-message transcripts make the
initial positions independent uniform draws with replacement from
$N=131{,}072$ outer leaves. After $q$ signatures there are $d=64q$ draws. If
$U_d$ is the number of distinct positions, then
\begin{equation}
  \mathbb E[U_d]
  =N\left(1-\left(1-\frac1N\right)^d\right).
  \label{eq:occupancy}
\end{equation}
Solving $\mathbb E[U_d]=32{,}768$ gives approximately 37,706.9 draws, or
589.17 signatures. This is an expectation, not a stopping guarantee. The
attack counts actual positions and continues until the threshold is reached.

Changing one draw changes $U_d$ by at most one. McDiarmid's bounded-difference
inequality therefore gives, for $\mu_d=\mathbb E[U_d]>32{,}768$,
\begin{equation}
 \Pr[U_d<32{,}768]
 \leq \exp\left(-\frac{2(\mu_d-32{,}768)^2}{d}\right).
 \label{eq:mcdiarmid}
\end{equation}
Table~\ref{tab:probability} instantiates the expectation and this conservative
bound. The observed counts closely match the occupancy expectation.

\begin{table}[ht]
\centering
\small
\begin{tabular}{@{}rrrrr@{}}
\toprule
$q$ & draws $d$ & $\mathbb E[U_d]$ & observed $U_d$ & failure bound \\
\midrule
641 & 41,024 & 35,224.60 & 35,203 & $<2^{-424}$ \\
650 & 41,600 & 35,644.88 & 35,603 & $<2^{-574}$ \\
1,023 & 65,472 & 51,534.14 & --- & $<2^{-15{,}520}$ \\
\bottomrule
\end{tabular}
\caption{Random-oracle occupancy estimates and measured distinct positions.
The bounds are analytic and do not replace the attack's explicit threshold
check.}
\label{tab:probability}
\end{table}

Both retained attacks chose a comfortable fixed query count; neither searched
for or claims the minimum successful count. Conditional on the random-oracle
model, an incremental adversary that is allowed the full 1,023-query budget
fails to collect enough positions with probability below the final row's
bound.

Once $k$ positions are available, standard product-tree interpolation costs
quasilinear to $O(k\log^2 k)$ field operations per row, depending on the
polynomial-arithmetic implementation. Four rows require $O(k)$ core state plus
temporary interpolation storage. The recovered $N$-element witness occupies
4 MiB at 32 bytes per extension-field element; the full encoded slab occupies
16 MiB, and the initial binary Merkle tree uses about 8 MiB of digests. Oracle
transcripts dominate storage: the independent 650-response corpus contained
219,867,416 serialized signature bytes. The measured peak resident memory
remained below 151 MiB during public recovery.

\section{Full-parameter experiments}
\label{sec:experiments}

\subsection{Primary implementation}

The primary experiment, \texttt{JEVIL-EXP-001}, used Rust 1.88.0 and the
pinned upstream archive. The oracle generator wrote only a serialized public
key, distinct query messages, ordinary signatures, range metadata, and hashes
to the recovery boundary. The public recovery executable had no input path or
API for the target secret seed or original signer cache. Every response passed
ordinary verification before extraction.

The attack tree left the ordinary signer and verifier files byte-identical to
the pinned source. Their SHA-256 digests were respectively
\begin{center}
\hash{34a61da541e7239402e2e49fedee76aaa50fc4c328c49cc9a307dca3bac2fb74}
\end{center}
and
\begin{center}
\hash{d84341b99aa3b76413c20118f7fb0a78886754e1d9c7c0f952c7f930037e569b}.
\end{center}
One preflight response and four disjoint 160-response shards covered indices 0
through 640. The public process received the five sealed public directories,
verified 641 signatures, interpolated the state, and forged independently.

\subsection{Hostile adversarial review}

A separately tasked reviewer checked query accounting, freshness, public-only
data flow, source identity, row semantics, root reconstruction, and the
ordinary-verifier boundary. A checker linked only to pristine pinned source
accepted the primary forgery and rejected changed-message, changed-root, and
altered-signature controls. The reviewer then replayed all five sealed public
input directories with no new signing query. It reproduced the 35,203
positions, row hashes, root, and controls, and generated a different accepted
fresh signature. The review verdict was
\texttt{PASS\_WITH\_LIMITATIONS / NO\_TECHNICAL\_VETO}.

\subsection{Fresh-context independent reproduction}

The \texttt{JEVIL-EXP-002} design was frozen before execution. A fresh-context
agent began from a new extraction of the pinned source and the paper-level
width-four observation. It separately wrote the observer, oracle, recovery,
interpolator, state adapter, and clean-verifier driver. It did not import the
primary implementation, key, transcript, openings, recovered rows, or public
artifacts. A new key seed was drawn from \texttt{/dev/urandom}; the recovery
process could not read the private seed directory.

The independent interpolator used a subproduct tree, a complementary erasure
locator derived from the product derivative, inverse radix-two transforms, and
fast reversed power-series division. A synthetic four-row test passed before
oracle collection. After recovery, the forgery was passed to a second source
extraction whose upstream files were untouched; only an external verifier
driver was added. Both the attack build and this untouched-source verifier
accepted.

\subsection{Results}

Table~\ref{tab:results} gives the two full-parameter trials. They used different
keys, disjoint message namespaces, independently written attack code, and
separate fresh proof randomness.

\begin{table}[ht]
\centering
\small
\begin{tabular}{@{}lrrrrl@{}}
\toprule
Run & queries & openings & distinct & fresh sig. bytes & verifier \\
\midrule
Primary & 641 & 41,024 & 35,203 & 343,532 & accept \\
Independent & 650 & 41,600 & 35,603 & 332,140 & accept \\
\bottomrule
\end{tabular}
\caption{Full-parameter state recovery and fresh-message forgery.}
\label{tab:results}
\end{table}

The primary result recovered all four rows, reproduced the target root and
$w=f(z)$, and signed an unqueried message. Its fresh signature SHA-256 was
\hash{3096cea57f85088a8a6c3eeba412b97a04a70ce940fd95696ff3b9f4b0f3f893}.
The independent run reached the same endpoint on a new key; its fresh signature
SHA-256 was
\hash{7e7b8fa3a213fcbfe6fe26792d5f8233df5268eb54914d217f3dd500e1066ed1}.
Neither fresh signature appeared among its oracle responses.

The retained full-parameter empirical result is two successes in two trials,
plus one public-data replay on the primary key. This does not justify a broad
empirical rate estimate. Equation~\eqref{eq:mcdiarmid} supplies the conditional
random-oracle success bound; every concrete run also checked the threshold
directly before interpolation.

\begin{table}[ht]
\centering
\small
\begin{tabular}{@{}lrrr@{}}
\toprule
Stage & Wall time & Peak RSS & workers \\
\midrule
Primary collection, each of four 160-query shards & 198.48--203.05 s & 131,060 KiB max & 1 each \\
Primary public recovery, signing, and controls & 486.90 s & 143,756 KiB & 1 \\
Hostile public-data replay & 475.63 s & 143,920 KiB & 1 \\
Independent 650-response collection & 13:19.13 & 131,256 KiB & 1 \\
Independent public recovery, signing, and controls & 4:13.12 & 150,380 KiB & 1 \\
Independent untouched-source verifier and controls & 0.93 s & 91,840 KiB & 1 \\
\bottomrule
\end{tabular}
\caption{Measured full-parameter resources. The primary collection shards were
independent processes; the other listed stages used one process and one logical
CPU.}
\label{tab:resources}
\end{table}

The experiments ran on x86-64 Linux with an AMD EPYC processor, 64 logical
CPUs, and 131,880,760 KiB installed memory. Each experiment process stayed
below a 4 GiB virtual-memory cap; no process approached the cap.

\subsection{Negative controls}

Both primary and independent runs enforced the following controls:
\begin{itemize}
  \item every oracle response passed the ordinary verifier;
  \item the forged message was absent from the signing-query set;
  \item the forged signature differed from every oracle response;
  \item changing the message caused rejection;
  \item changing the public root caused rejection;
  \item altering the serialized signature caused rejection or parse failure;
  \item corrupting a recovered row failed root or public-evaluation binding;
        and
  \item signature 1,024 and zero-based query index 1,023 were never requested
        or consumed.
\end{itemize}
These controls rule out replay, a weakened attack verifier, wrong-root
acceptance, coefficient recovery without a signing consequence, and accidental
use of the advertised cliff.

\section{Why this is below the advertised cliff}
\label{sec:cliff}

The advertised mechanism and this attack use different public information.
Jevil's stated cliff counts the $K=16$ message-derived evaluations
$(x_t,f(x_t))$ in each signature and sets the polynomial dimension to
$(n^\star+1)K$. Its trigger is signature number 1,024.

Our attack ignores that recovery route. It accumulates the 64 initial
width-four zk-WHIR leaves that ordinary proofs must reveal. The primary attack
stopped at signature number 641; the independent reproduction stopped at 650.
Their largest zero-based query indices were 640 and 649. Both are within the
claimed operating budget and hundreds of signatures before the designed
trigger.

The distinction also explains why recovering $f$ alone would be insufficient.
Ordinary signing needs the commitment state, including $r_{\rm zk}$, to open the
existing public root. The width-four route recovers that complete state. Fresh
opening randomness then yields a reusable signer for the same public key.

\section{Prior public records and attack-history limits}
\label{sec:history}

The final exact-target search covered the exact title and URL, ePrint and
author indexes, citations, public repositories, issues, commits, advisories,
crates, and combinations of Jevil with attack, cryptanalysis, forgery,
width-four recovery, state recovery, and the measured query counts. It found no
prior public practical signing attack against the frozen current profile within
1,023 ordinary queries. Two public records remain material and are disclosed
here.

\paragraph{26 May root-binding forgery fix.}
Author-repository commit
\hash{3c7968ddacf4bddcbd178b3e7a61625148dabf9c} bound WHIR verification to
\path{pk.root} and added a regression named
\path{forged_root_attack_now_rejected}~\cite{rootfix}. It documents a
forgery against an earlier development state. The fix predates the paper's
29 May 2026 receipt and is present in the frozen target; the old forged-root
proof does not verify under the current ordinary verifier. This public record
precludes any claim that no Jevil development version had a forgery.

\paragraph{2 June non-power-of-two Merkle-helper report.}
A Swival trace published in the \emph{security-audits} repository at commit
\hash{cea51fd3cec15036e4d845d0fb6d3f5628709d60} shows the multiproof verifier
accepting padded index 13 for an artificial 13-leaf tree~\cite{merklebug}. Its
embedded \texttt{src/merkle.rs} is byte-identical to
the frozen file. The audited Jevil profile constructs the relevant commitment
sizes as powers of two, and ordinary verification derives positions modulo the
codeword length instead of accepting an attacker-supplied padded index. The
trace demonstrates no signing-state recovery, fresh-message forgery, or
width-four accumulation attack against this profile. It is nevertheless a
previously reported defect in the current source file and should not be
omitted.

The history conclusion is therefore deliberately narrow: this is a practical
below-budget attack on the current fourth-revision, source-backed profile. We
make no unqualified priority statement.

\section{Reproducibility and evidence identity}
\label{sec:reproducibility}

The primary and independent artifacts retain prospective designs, source
manifests, locked dependency graphs, exact commands, machine descriptions,
public query corpora, logs, results, and negative-control outputs. The critical
sealed identities appear in Table~\ref{tab:seals}.

\begin{table}[ht]
\centering
\small
\begin{tabularx}{\textwidth}{@{}lY@{}}
\toprule
Object & SHA-256 \\
\midrule
Primary result record &
\hash{36b01f7df957752489a163eabcce200c38e269f8a34f50b6045fe68c7591f3d4} \\
Primary 1,350-entry evidence manifest &
\hash{4f1b226bbdb82987bd3388d85e90e94362a5d5ceec397c7a2a062e2f1534f5b9} \\
Hostile-review handoff &
\hash{ad8af944d0b6b049f683c30158b43957a017b97d4bdcf72adbf325992018b8b6} \\
Independent prospective design &
\hash{da0fd7bd14121b048ae8bdfbf338ac2afa930de60e14474f2a71e18e7f08dd8d} \\
Independent result record &
\hash{0d43f5ed095de93a12e16bfb96a2a6737654727b2488a9eb4431d5fe20793ce8} \\
Independent 1,453-entry manifest &
\hash{02414db4b8bd5d12d0065ae1fd452c46f07705b5380e769bb3b70018909a1e4c} \\
\bottomrule
\end{tabularx}
\caption{Evidence seals for the reported results.}
\label{tab:seals}
\end{table}

To reproduce the endpoint, a run must pin the target above, generate a fresh
$n^\star=1023$ key, query no more than 1,023 distinct ordinary messages, and
give its recovery process only the public key, messages, and serialized
ordinary signatures. It must authenticate and verify openings before use,
recover all four rows, match $w=f(z)$ and the exact root, sign an unqueried
message with fresh randomness, and obtain acceptance from an ordinary verifier
whose upstream acceptance logic is unchanged. A result that only recovers
coefficients, substitutes an attack-specific verifier, or consumes signature
1,024 does not reproduce this claim.

\section{Limitations and design implications}

The following limits are material to interpreting the result.
\begin{itemize}
  \item The executable attack targets the frozen source commit and its
        $n^\star=1023,c=4$ row. The source predates the final PDF, revisions
        1--3 were unavailable, and the final PDF does not pin a commit.
  \item The intended Definition~3.22 four-row result is analytic. The full
        attack was executed against the official source layout, whose
        contiguous split is nonconforming. The full-recovery threshold is the
        same, but earlier partial leakage differs. The paper's literal scalar
        stripe map is under-specified.
  \item The primary run and one independent fresh-key run are two successful
        full-parameter trials, not a many-key empirical study. The probability
        bounds assume the paper's Fiat--Shamir random-oracle model.
  \item Process isolation was enforced by separate executables, public-only
        argument and file boundaries, source audits, a scrubbed environment,
        and an unreadable independent seed directory. A mount-namespace
        sandbox and syscall trace were unavailable under the managed kernel
        policy.
  \item The primary key used a deterministic seed retained in the separate
        oracle-generator source, and the primary seal did not hash the exact
        primary-time release executable. Static data-flow and binary checks
        found no recovery path to that seed, the hostile reviewer replayed the
        sealed public inputs, and the independent run used a fresh random key
        whose seed was unavailable to its recovery process.
  \item Public attack-history search cannot exclude private, deleted,
        unpublished, or unindexed work, and several bibliographic endpoints
        were challenged or rate-limited.
\end{itemize}

The immediate design lesson is to account for one width-four leaf as four
scalar disclosures, or as one disclosure in each of four separately masked
rows. For the stated Definition~3.22 padding, the multi-opening hiding
condition is $q\theta\leq B/4$, which permits at most 448 signatures at
$\theta=64$, not 1,023. Merely lowering a counter avoids this demonstrated
full-state threshold only by changing Jevil's advertised row. Preserving the
1,023-signature goal requires redesign: for example, enlarge the per-row
dimension and masking budget consistently, avoid exposing reusable raw initial
leaves, or refresh the commitment in a way that remains bound to the public
key. Each option changes parameters or protocol structure and needs a new
security analysis.

\section{Conclusion}

Jevil's ordinary proof reveals four initial Reed--Solomon values at each of 64
positions per signature. Because the commitment is reused and each of four
rows has dimension 32,768, ordinary signatures accumulate enough authenticated
positions to recover the entire committed state well before the intended
1,024th-signature cliff. The primary implementation did so with 641 signing
queries; an independently written fresh-key implementation did so with 650.
Both rebuilt the exact public state and produced fresh-message signatures
accepted by untouched ordinary verification logic. The attack is standard
interpolation applied to a width and budget mismatch, but its signing
consequence is complete: it yields an equivalent signer inside the claimed
1,023-signature operating range.

\section*{Acknowledgements}
The computations were performed on the Norwegian Research and Education Cloud (NREC), using resources provided by the University of Bergen and the University of Oslo.

\begin{thebibliography}{9}

\bibitem{jevil}
N. Kobeissi,
``Jevil: A Catastrophic-Failure-by-Design Signature Scheme,''
IACR Cryptology ePrint Archive, Report 2026/1103, revision 4, 2026.
\url{https://eprint.iacr.org/2026/1103}

\bibitem{zkwhir}
A. Chiesa, G. Fenzi, and G. Weissenberg,
``Zero-Knowledge IOPPs for Constrained Interleaved Codes,''
IACR Cryptology ePrint Archive, Report 2026/391, 2026.
\url{https://eprint.iacr.org/2026/391}

\bibitem{reedsolomon}
I. S. Reed and G. Solomon,
``Polynomial Codes over Certain Finite Fields,''
\emph{Journal of the Society for Industrial and Applied Mathematics},
vol.~8, no.~2, pp.~300--304, 1960.

\bibitem{mca}
J. von zur Gathen and J. Gerhard,
\emph{Modern Computer Algebra}, 3rd ed., Cambridge University Press, 2013.

\bibitem{rootfix}
N. Kobeissi,
``Bind WHIR proof verification to the public-key root,''
\texttt{symbolicsoft/jevil} commit
\hash{3c7968ddacf4bddcbd178b3e7a61625148dabf9c}, 26 May 2026.
\url{https://github.com/symbolicsoft/jevil/commit/3c7968ddacf4bddcbd178b3e7a61625148dabf9c}

\bibitem{merklebug}
\texttt{jedisct1/security-audits},
``multiproof verifier accepts out-of-range indices,'' Swival trace in commit
\hash{cea51fd3cec15036e4d845d0fb6d3f5628709d60}, 2 June 2026.
\url{https://huggingface.co/datasets/jedisct1/security-audits/commit/cea51fd3cec15036e4d845d0fb6d3f5628709d60}

\end{thebibliography}

\end{document}
